\documentclass[11pt,aspectratio=169]{beamer}

\usepackage{listings}


\usepackage{hyperref}
\usepackage{tikz}
\usepackage{calc}

\hypersetup{
  urlcolor=gray,
  colorlinks=true
  }

\setbeamertemplate{navigation symbols}{}
\setbeamertemplate{footline}[frame number]

\setbeamercolor{itemize item}{fg=blue}
\setbeamercolor{itemize subitem}{fg=blue}
\setbeamercolor{enumerate item}{fg=blue}
\setbeamercolor{enumerate subitem}{fg=blue}

\setbeamercolor{title}{fg=black,bg=gray!50!white}
\setbeamercolor*{subtitle}{fg=white}
\setbeamercolor*{date}{fg=gray,bg=black}
\setbeamercolor{author}{fg=gray,bg=black}
\setbeamercolor{institute}{fg=gray,bg=black}

\newcommand{\nrcanTitleframe}{
{\usebackgroundtemplate{%
    \includegraphics[width=\paperwidth,height=\paperheight]{nrcan_title}
  }
  \begin{frame}[plain]
  \titlepage
\end{frame}
}}

\newenvironment{nrcanFrame}{
 \setbeamertemplate{footline}[frame number]{}    
\usebackgroundtemplate{%
    \includegraphics[width=\paperwidth,height=\paperheight]{nrcan_content}}
  \begin{frame}
}{
  \end{frame}
}

\newcommand{\nrcanLastFrame}{
\setbeamertemplate{footline}[frame number]{}
  \usebackgroundtemplate{%
    \includegraphics[width=\paperwidth,height=\paperheight]{nrcan_lastpage}
  }
  \begin{frame}
    \begin{tikzpicture}
      \node at (0,0) {};
      \node[scale=.7] at (9cm,-7.5cm) {\tiny Made with: \url{https://github.com/gnewton/nrcan_latex_template}};
    \end{tikzpicture}
  \end{frame}
}

%%%%
\newcommand{\nrcanGraphicFrame}[4]{
  % 1 title
  % 2 graphic
  % 3 width/height
  % 4 url source
  \def\Filename{#2} % From https://tex.stackexchange.com/questions/474504/using-optional-argument-in-closing-statement-of-new-environment
  \usebackgroundtemplate{}
    \setbeamertemplate{footline}{\centering{\hspace{2mm}\tiny Source: \url{#4}}
\hfill%
  \usebeamercolor[fg]{page number in head/foot}%
  \usebeamerfont{page number in head/foot}%
  \insertframenumber\,/\,\inserttotalframenumber\kern1em\vskip2pt%
}

  \begin{frame}
   
    \frametitle{#1}
    \centering
        \includegraphics[#3]{#2}
  \end{frame}
}
%%%%
% Change the templace for the title page     
\defbeamertemplate*{title page}{customized}[1][]
{
  \begin{tikzpicture}
    \node[rounded corners,align=center,inner sep=5mm,text width=13cm,text=black,fill=gray!50!white,opacity=.8, text opacity=1] at (0,0)
         {%
           {\huge \inserttitle} \\
           {\Large \insertsubtitle} \ \\
           \ \\
               {\Large \insertauthor}\\
               \ \\
                   {\small \insertinstitute}
                   \ \\
      \ \\
      \insertdate};
\end{tikzpicture}
}

%%%%%%%%%%



\title[Crisis] % (optional, only for long titles)
{Microservices: APIs, Lambdas and FAAS}
\author[Newton] % (optional, for multiple authors)
{$\lambda$ Glen ~Newton $\lambda$ }
\institute[Dept This and That] % (optional)
{
  Section Head, Scientific Computing\\
  Geological Survey of Canada\\
  Natural Resources Canada
}
\date[2021 10 28] % (optional)
     {NRCan Sectors Cloud Working Group\\
     2021 10 28}

\logo{\begin{tikzpicture}
    \node at (0cm,0cm) {};
        \node at (0cm,18cm) {};
 \filldraw[color=blue!50, fill=blue!25, very thick](1mm,17.4cm) circle (0.5);
 \node[color=white] at (1mm,17.4cm) {\Huge $\lambda$};
 \end{tikzpicture}}



%%%%%%%%%%%%%%%%%%%%%%%%%%%%%%%%%%%%%%%%%%%%%%%%%%%%%%%%%%%%%%%%%%%%%%%%%%%%%%%
%%%%%%%%%%%%%%%%%%%%%%%%%%%%%%%%%%%%%%%%%%%%%%%%%%%%%%%%%%%%%%%%%%%%%%%%%%%%%%%
\begin{document}


%%%%%%%%%%%%%%%%%%%%%%%%%%%%%%%%%%%%%%%%%%%%%%%%%%%%%%%%%%%%%%%%%%%%%%%%%%%%%%%
%% Title page
\nrcanTitleframe

  
%%%%%%%%%%%%%%%%%%%%%%%%%%%%%%%%%%%%%%%%%%%%%%%%%%%%%%%%%%%%%%%%%%%%%%%%%%%%%%%
\begin{nrcanFrame}
  \frametitle{Introduction}
  \begin{itemize}
  \item FAAS (Function-as-a-Service) in the \textit{X}-as-a-service universe
  \item What is FAAS and how does it work?
  \item How useful is it / what are the use cases?
  \item Deep dive using AWS Lambda
  \item What are the advantages/disadvantages?

    \end{itemize}
\end{nrcanFrame}

\begin{nrcanFrame}
  \frametitle{What is FAAS? Wikipedia: FAAS}
  \textit{``Function as a service (FaaS) is a category of cloud
    computing services that provides a platform allowing customers to
    develop, run, and manage application functionalities \alert{without the
    complexity of building and maintaining the infrastructure
    typically associated with developing and launching an app}.
    Building an application following this model is one way of
    achieving a "serverless" architecture, and is typically used when
    building \alert{microservices} applications.  
\\
      \vspace*{2mm}
      FaaS was initially offered by various start-ups circa 2010, such
      as PiCloud. 
      \\
\vspace{2mm}
AWS \alert{Lambda} was the first FaaS offering by a large public cloud
vendor, followed by Google \alert{Cloud Functions}, Microsoft \alert{Azure Functions},
IBM/Apache's OpenWhisk (open source) in 2016 and Oracle Cloud Fn (open
source) in 2017.'' 
}

  {\tiny Source: \url{https://en.wikipedia.org/wiki/Function_as_a_service}}
\end{nrcanFrame}
    
\nrcanGraphicFrame
  {Car as a Service}
  {carServiceAnalogy}
  {height=\paperheight/5*4,width=\paperwidth/7*6}
  {https://rajeevkuruganti.com/2019/cloud-iaas-paas-saas-and-more/}

\nrcanGraphicFrame
  {Hosting Models: Traditional, Cloud, Cloud Native}
  {cloudStack0}
  {width=\paperwidth/7*6,height=\paperheight/8*7}
  {https://devops.stackexchange.com/questions/5688/how-can-caas-paas-and-faas-users-know-if-the-operating-system-of-their-server}

\nrcanGraphicFrame
  {Customer-managed unit of scale}
  {cloudStacks}
  {width=\paperwidth/4*3}
  {https://serverless.zone/abstracting-the-back-end-with-faas-e5e80e837362}

\begin{nrcanFrame}
  \frametitle{FAAS Characteristics 1/2}
  \begin{itemize}
  \item Serverless: server completely abstracted / hidden from
    customer: No hardware, server, VM, OS, etc to manage
  \item Focus on code rather than infrastructure
  \item Basically short-lived functions that do a specific task
  \item Pay only for use / run only when needed
  \item Event driven
  \end{itemize}
\end{nrcanFrame}


\begin{nrcanFrame}
  \frametitle{FAAS Characteristics 2/2}
  \begin{itemize}
  \item Scalable by design
  \item Concurrent by design
  \item Reliable / redundant 
  \item Provisioning/deployment (usually) simple
  \item \alert{Secure}: Functions (should) have IAM roles that 
    allow them to do exactly
    alert{(and only)} what they need to do.
  \end{itemize}
\end{nrcanFrame}



\begin{nrcanFrame}
  \frametitle{FAAS Characteristics: Software developer best practices}
    \begin{itemize}
  \item Functions (should) perform only one action
  \item Small, well-defined and specific scope
  \item Isolation of concerns
  \item Stateless (with exceptions)
  \item Limit dependencies / reduce deployment package size. For some languages, the loading of
    libraries will add to latency
  \item Within functions, separate cloud vendor-specific handler code from function business
    logic.
    Allows for unit testing and (somewhat) limits vendor lock-in.
    \end{itemize}
\end{nrcanFrame}




\begin{nrcanFrame}
  \frametitle{FAAS/Serverless: Best Practices: Some references}
  \begin{itemize}
  \item AWS
    \textit{\href{https://docs.aws.amazon.com/lambda/latest/dg/best-practices.html}
      {Best practices for working with AWS Lambda functions}}

  \item AWS
    \textit{\href{https://www.intentsg.com/aws-lambda-serverless-coding-best-practices/}
      {AWS Lambda Serverless Coding Best Practices}}

  \item Azure
    \textit{\href{https://docs.microsoft.com/en-us/azure/azure-functions/functions-best-practices?tabs=csharp}
      {Best practices for reliable Azure Functions}}

  \item Azure
    \textit{\href{https://www.devbridge.com/articles/how-when-and-why-to-use-microsoft-azure-cloud-services/}
      {How, when, and why to use Microsoft Azure Cloud Services Azure Functions}}

    \item
      \textit{\href{https://martinfowler.com/articles/serverless.html}
        {Serverless Architectures}}
      
  \item
    \textit{\href{https://www.captechconsulting.com/blogs/applying-microservice-patterns-best-practices-to-faas-part-1-faas-overview}
      {Applying Microservice Patterns \& Best Practices To FaaS (Part 1 - FaaS Overview)}}
    \end{itemize}
\end{nrcanFrame}



\begin{nrcanFrame}
  \frametitle{AWS Lambda: Case Study}
  \begin{itemize}
  \item AWS Lambda  = AWS FAAS
  \item Supports multiple languages: Java, Go, PowerShell, Node.js, C\#, Python,Ruby + \href{https://docs.aws.amazon.com/lambda/latest/dg/runtimes-custom.html}{custom runtimes} (any programming language)
  \item Packaging:
    \href{https://docs.aws.amazon.com/lambda/latest/dg/gettingstarted-package.html}
    {zip file or container}
  \item Event driven
  \item Pay per use (execution time * memory): \href{https://aws.amazon.com/blogs/aws/new-for-aws-lambda-1ms-billing-granularity-adds-cost-savings/}{1ms billing granularity}
  \end{itemize}
\end{nrcanFrame}


\begin{nrcanFrame}
  \frametitle{AWS Lambda: Default limits 1/2}
  \begin{itemize}
  \item Local \textit{ephemeral} storage \texttt{/tmp}: 500 MB 
  \item Max concurrency: 1000 instances / region
  \item Memory: 128 MB to 10 GB
  \item CPU: Proportional to memory (see below)
  \item Mx running time (timeout): 15min
  \end{itemize}
\end{nrcanFrame}




\begin{nrcanFrame}
  \frametitle{AWS Lambda: Default limits 2/2}
  \begin{itemize}
  \item Invocation payload: 6 MB (synchronous); 256 KB (asynchronous)
  \item Deployment package (zip) size: 50 MB (zipped), 250 MB (unzipped)
  \item Container image code package size limit: 10 GB
  \item (As of June 2020): Lambdas can mount \href{https://aws.amazon.com/efs/}{EFS (elastic file system)}
    \begin{itemize}
    \item This allows lambdas to be stateful across invocations and lambdas, with state stored in EFS volume
    \item Mounting adds to cold start latency (see below):
      \href{https://aws.amazon.com/blogs/compute/using-amazon-efs-for-aws-lambda-in-your-serverless-applications/}{``hundreds
        of milliseconds''} of latency
    \item If using EFS, lambda needs to be in the same
      \href{https://aws.amazon.com/vpc/}{VPC (virtual private cloud)} as the EFS volume
    \end{itemize}
  \end{itemize}
\end{nrcanFrame}


\begin{nrcanFrame}
  \frametitle{AWS Lambda: Proportion of vCPU to memory}
  \centering
  \begin{table}
  \begin{tabular}{|c||l|}
    \hline    \hline
    \bf Memory & \# \bf vCPU \\  \hline
    128 MB & 1 vCPU \\ \hline
  832 MB & 2 vCPU \\ \hline
  3 GB & 3 vCPU \\ \hline
  5.3 GB & 4 vCPU \\ \hline
  7 GB & 5 vCPU \\ \hline
  8.8+ GB & 6 vCPU \\
      \hline    \hline
  \end{tabular}
  \end{table}
\end{nrcanFrame}



\begin{nrcanFrame}[fragile=singleslide]
    \frametitle{Go language lambda boilerplate}
  \centering
{\scriptsize
  \begin{lstlisting}[language=go]
package main

import (
        "fmt"
        "context"
        "github.com/aws/aws-lambda-go/lambda"
)

type MyEvent struct {
        Name string `json:"name"`
}

func HandleRequest(ctx context.Context, name MyEvent) (string, error) {
        return fmt.Sprintf("Hello %s!", name.Name ), nil
}

func main() {
        lambda.Start(HandleRequest)
}
  \end{lstlisting}
}
  {\tiny Source: \url{https://docs.aws.amazon.com/lambda/latest/dg/golang-handler.html}}
\end{nrcanFrame}





  \begin{nrcanFrame}[fragile=singleslide]
    \frametitle{Python language lambda boilerplate}
               {\footnotesize
                 \begin{lstlisting}[language=python]
def lambda_handler(event, context):
    message = 'Hello {} {}!'.format(event['first_name'], event['last_name'])  
    return { 
        'message' : message
    }
                 \end{lstlisting}
                 
                 \vspace{3cm}
                        {\tiny Source: \url{https://docs.aws.amazon.com/lambda/latest/dg/python-handler.html}}}
  \end{nrcanFrame}


\nrcanGraphicFrame
  {Example architecture showing both HTTP and S3 events}
  {awsLambdaEventArch}
  {height=\paperheight*4/5}
  {https://programmaticponderings.com/2019/10/04/event-driven-serverless-architectures-with-aws-lambda-sqs-dynamodb-and-api-gateway/}


\begin{nrcanFrame}
  \frametitle{Lambda containers}
  \begin{itemize}
  \item Supported containers: Docker or Open Container Initiative (OCI)
  \item Deploy from \href{https://aws.amazon.com/ecr/}{Amazon Elastic Container Registry (ECR)}
  \item Many AWS-supplied base image options with \href{https://docs.aws.amazon.com/lambda/latest/dg/lambda-runtimes.html}{pre-installed runtimes}
  \item
    AWS - \href{https://docs.aws.amazon.com/lambda/latest/dg/images-create.html}{Creating
      Lambda container images}
    \item Azure Functions \textit{do} support Docker (but only for certain
    \href{https://docs.microsoft.com/en-us/azure/azure-functions/functions-scale}{hosting
      plans})
    \item Google Cloud has \href{https://cloud.google.com/run}{Cloud
      Run}, basically like AWS Lambda containers
  \end{itemize}
\end{nrcanFrame}


\begin{nrcanFrame}
  \frametitle{Cold start/latency}
  When the Lambda is invoked, the following steps take place, each taking time and adding to the latency of the function call:
  \begin{enumerate}
  \item Code download (zip from S3 or container from ECR)
  \item Start execution environment
  \item Run initialization code
  \item Run handler code
  \end{enumerate}

  \ \\
  $\#1 + \#2$ are the \alert{cold start}.\\
  \textit{Cold start} can be avoided using \href{https://docs.aws.amazon.com/lambda/latest/dg/provisioned-concurrency.html}{provisioned concurrency}.\\
  Full explanation see \href{https://aws.amazon.com/blogs/compute/operating-lambda-performance-optimization-part-1/}{Operating Lambda: Performance optimization – Part 1}
\end{nrcanFrame}

\nrcanGraphicFrame
  {FAAS Cold Start Comparison: AWS, Azure, GCP}
  {coldStartComparison}
  {width=\paperwidth*10/11}
  {https://mikhail.io/serverless/coldstarts/big3/}

  \nrcanGraphicFrame
  {Cold Starts in Azure Functions: Windows vs. Linux}
  {azureWindowsVsLinux.png}
  {width=\paperwidth*10/11}
  {https://mikhail.io/serverless/coldstarts/azure/}

  \nrcanGraphicFrame
  {Cold Starts in Azure Functions: Impact of Package Size}
  {azurePackageSize.png}
  {width=\paperwidth*10/11}
  {https://mikhail.io/serverless/coldstarts/azure/}


\begin{nrcanFrame}
  \frametitle{Cold start/latency references}
  \begin{enumerate}
  \item GCP \textit{\href{https://www.infoq.com/news/2021/09/cloud-functions-min-instances/}{Google Introduces Minimum Instances to Reduce Cold Starts}}
  \item Comparison \textit{\href{https://mikhail.io/serverless/coldstarts/big3/}{Comparison of Cold Starts in Serverless Functions across AWS, Azure, and GCP}}
  \item Azure \href{https://mikhail.io/serverless/coldstarts/azure/}{Cold Starts in Azure Functions}
  \end{enumerate}

\end{nrcanFrame}


\begin{nrcanFrame}
  \frametitle{Lambda Lifecycle 1/2}
  The AWS lifecycle phases are:
  \begin{enumerate}
  \item Init: Happens at first invocation, or, with provisioned
    concurrency, in advance.
    \begin{enumerate}
    \item Extension init
    \item Runtime init
    \item Function init
    \end{enumerate}
  \item Invoke
  \item Shutdown 
    \begin{enumerate}
    \item Runtime shutdown
    \item Extension shutdown
    \end{enumerate}
  \end{enumerate}
\end{nrcanFrame}




\begin{nrcanFrame}
  \frametitle{Lambda Lifecycle 2/3}
  \begin{itemize}
    \item When a lambda \textit{instance} has responded to an event and completed, it is \alert{not immediately removed/deleted}:
  its environment is kept for an undefined period of time.
\item If the lambda is reinvoked during this time, this \textit{instance} may be used to respond to the request.
  \alert{If this happens, there is no cold start}.
\item This
  \href{https://xebia.com/blog/til-that-aws-lambda-terminates-instances-preemptively/}{experiment}
  from Sept 2020 suggests that lambdas are
  \textit{``\dots terminated
    if a  function isn’t invoked for 15 minutes \ldots didn’t find the
    ‘guaranteed to be cut-off’ time, but it’s somewhere between 10 and
    15 minutes of inactivity.''} -- Caveat emptor
\item 
  AWS
  \href{https://docs.aws.amazon.com/lambda/latest/dg/runtimes-context.html\#runtimes-lifecycle-shutdown}{suggests}
  that lambda instances can \textit{cache} data in their \texttt{/tmp} directories, just in case they are reinvoked.
  Additional function logic is needed to take advantage of this situation.
  If provisioned concurrency is used, then this strategy is a more
  viable option.
  \end{itemize}
\end{nrcanFrame}

\nrcanGraphicFrame
  {Lambda Lifecycle 3/3}
  {OverviewSuccessfulInvokes}
  {width=\paperwidth*10/11}
  {https://docs.aws.amazon.com/lambda/latest/dg/runtimes-context.html}



\begin{nrcanFrame}
  \frametitle{Cold start for different languages (ms), different memory}
  \centering
  \begin{table}
  \begin{tabular}{|r||r|r|r|r|r|r|r|r|}\hline    \hline
    \bf Memory      &\bf Java 	&\bf Graalvm& \bf .Net &	\bf Go& 	\bf Rust& 	\bf Python& 	\bf NodeJS& 	\bf Ruby\\ \hline
    128mb 	&OOM 	&1480 &	  11810 &	1050 &	844 &	641 &	        1190 &	773\\ \hline
    256mb 	&6570 &	774 &	  5820 &	661 &	480 &	527 &	        769 &	612\\ \hline
    512mb 	&5180 &	684 &	  2940 &	404 &	304 &	502 &	        771 &	677\\ \hline
    1024mb 	&4450 &	531 &	  1500 &	299 &	234 &	482 &	        656 &	652\\ \hline
    10240mb     &2790 &	501 &	  904 &	        327 &	219 &	449 &	        518 &	649\\ \hline
      \hline 
  \end{tabular}
  \end{table}

  {\tiny Source: \url{https://filia-aleks.medium.com/benchmarking-all-aws-lambda-runtimes-in-2021-cold-start-part-1-e4146fe89385}}
\end{nrcanFrame}



\nrcanGraphicFrame
  {128 MB Average Warm Start (ms)}
  {warmStart}
  {width=\paperwidth/7*6}
  {https://filia-aleks.medium.com/aws-lambda-battle-2021-performance-comparison-for-all-languages-c1b441005fd1}


\begin{nrcanFrame}
  \frametitle{Lambda cost exercise 1/4}
  Example \textit{modest} REST lambda microservice:
  \begin{itemize}
  \item 60 requests / minute ($\sim$86k req/day; $\sim$2.6m req/month)
  \item 512 MB lambda memory (1vCPU)
  \item 2 second (2000ms) average run time
  \item AWS Region: Canada (Central)
  \end{itemize}
\end{nrcanFrame}


\begin{nrcanFrame}
  \frametitle{Lambda cost exercise 2/4: Lambda \textit{only} costs, per month}
  Using
  \href{https://calculator.aws/\#/createCalculator/Lambda}
       {AWS Pricing Calculator}:
  \ \\
      {\footnotesize
        \begin{itemize}
        \item \texttt{2,592,000 requests x 2,000 ms x 0.001 ms to sec conversion factor = 5,184,000.00 total compute (seconds)}
        \item \texttt{0.50 GB x 5,184,000.00 seconds = 2,592,000.00 total compute (GB-s)}
        \item \texttt{2,592,000.00 GB-s x 0.0000166667 USD = 43.20 USD (monthly compute charges)}
        \item \texttt{2,592,000 requests x 0.0000002 USD = 0.52 USD (monthly request) charges)}
        \end{itemize}
        }
            \ \\
      \texttt{\$43.20 USD + \$0.52 USD = \textbf{\$43.72 USD (\$54.14 CAD)}}

\end{nrcanFrame}



\begin{nrcanFrame}
  \frametitle{Lambda cost exercise 3/4: API Gateway costs}
  Using
  \href{https://calculator.aws/\#/createCalculator/APIGateway}{AWS
    Pricing Calculator}:
  \ \\ \vspace{1mm}
        {\footnotesize 
            \begin{itemize}
            \item
              \texttt{2.592 requests x 1,000,000 unit multiplier = 2,592,000 total REST API requests}
            \item
              \texttt{Tiered price for: 2592000 requests}
            \item \texttt{2592000 requests x 0.0000035000 USD = 9.07 USD}
            \item \texttt{Total tier cost = 9.0720 USD (REST API requests)}
            \item \texttt{Tiered price total for REST API requests: 9.072 USD}
            \item \texttt{0 USD per hour x 730 hours in a month = 0.00 USD for cache memory}
            \item \texttt{Dedicated cache memory total price: 0.00 USD}
            \end{itemize}
            }
    \ \\
    \texttt{API Gateway cost (monthly): \textbf{\$9.07 USD (\$11.23 CAD})}\\

\end{nrcanFrame}

        \begin{nrcanFrame}
  \frametitle{Lambda cost exercise 4/4: Total costs}
    \ \\ \vspace{1mm}
\texttt{\textbf{Total per month (Lambda + API Gateway): \$65.37 CAD}}
\ \\
\vspace{1cm}
    { \alert{NB: Does not cost-out backend AWS DB costs: most applications
        would likely have this additional cost. But this would be the same for all
        implementations, i.e. Lambda,
        \href{https://aws.amazon.com/ec2/}{EC2},
        \href{https://aws.amazon.com/fargate/?c=ser&sec=srv}{Fargate},
        \href{https://aws.amazon.com/eks/?pg=ln&sec=hiw}{Elastic Kubernetes
          Service (EKS)} etc\ldots
      }
    }
        \end{nrcanFrame}

\begin{nrcanFrame}
  \frametitle{Azure, GCP costing differences}
  \begin{itemize}
  \item Azure: billing precision: 1s; Monthly subscriptions
  \item GCP: billing precision: 100 ms
    \end{itemize}
  \end{nrcanFrame}

\nrcanGraphicFrame
  {FAAS Decision Tree (?)}
  {faas-decision-tree}
  {width=\paperwidth/7*6}
  {https://developer.ibm.com/articles/when-to-use-iaas-faas-paas-and-caas/}


\begin{nrcanFrame}
  \frametitle{FAAS Disadvantages (1/2)}
    \begin{itemize}
    \item Needs tuning (to optimize cost vs. performance): Memory and \# of vCPUs.
      Tools exist to automate this: AWS:
      \href{https://docs.aws.amazon.com/lambda/latest/operatorguide/profile-functions.html}
           {AWS Lambda Power Tuning}. \alert{Minor}.
  \item Limits on memory and \# vCPUs: max 10 GB RAM and/or 6 vCPUs (AWS
    Lambda): Does not support many scientific use cases (i.e. large
    memory work loads, etc).
  \item Limits on running time (15min for AWS): Too short for some enterprise and scientific workloads.
    For some use cases, this can be overcome with \href{https://en.wikipedia.org/wiki/Application\_checkpointing}{checkpointing} (often used in traditional high performance computing HPC) and lambda \href{https://en.wikipedia.org/wiki/Pipeline\_(software)}{pipelining}.
  \item Cold start issues
    \end{itemize}
\end{nrcanFrame}

\begin{nrcanFrame}
  \frametitle{FAAS Disadvantages (2/2)}
    \begin{itemize}
  \item Non-traditional architecture and technologies can challenge
    some developers / architects
  \item Stateless (with exceptions), which can be a challenge to designing and architecting
  \item Vendor lock-in: Not too high a risk: lock-in is more likely
    with the use of vendor BAAS (backend as a service) services like
    backend databases, that the lambda uses
  \item Basic monitoring and debugging : not as mature as other stacks,
    can be problematic.
    Improving. See: AWS: \href{https://docs.aws.amazon.com/lambda/latest/operatorguide/monitoring-observability.html}{Monitoring and observability}; Azure: \href{https://docs.microsoft.com/en-us/azure/azure-functions/functions-monitoring}{Monitor Azure Functions}
    \end{itemize}
\end{nrcanFrame}


\begin{nrcanFrame}
  \frametitle{FAAS Orchestration}
  ``\textit{Frequently, orchestration is what we actually mean when we
    are talking about automating. Orchestration is automating many
    tasks together. It’s automation not of a single task but an entire
    IT-driven process. Orchestrating a process, then, is automating a
    series of individual tasks to work together.}
  -\href{https://www.bmc.com/blogs/it-orchestration-vs-automation-whats-the-difference/}
       {BMC Blogs, Stephen Watts, Sept 2020}

\end{nrcanFrame}

\begin{nrcanFrame}
  \frametitle{FAAS Orchestration: AWS}
  \begin{itemize}
  \item AWS: Step Functions -
             \href{https://docs.aws.amazon.com/step-functions/index.html}{\textit{``Step Functions is a serverless
      orchestration service that lets you combine AWS Lambda functions
      and other AWS services to build business-critical
      applications}...\href{https://aws.amazon.com/step-functions/?step-functions.sort-by=item.additionalFields.postDateTime&step-functions.sort-order=desc}{
           Workflows \alert{manage failures, retries, parallelization,}
           service integrations, and observability so developers can
           focus on higher-value business logic.}}''
    \begin{itemize}
    \item State machine, encoded into JSON
    \item API to run \alert{and build}
    \item Visual editor
      \href{https://docs.aws.amazon.com/step-functions/latest/dg/workflow-studio.html}
      {Workflow Studio}
    \end{itemize}
  \end{itemize}
\end{nrcanFrame}

\begin{nrcanFrame}
  \frametitle{FAAS Orchestration: Azure}
    \begin{itemize}
    \item \href{https://docs.microsoft.com/en-us/azure/azure-functions/durable/}
         {Durable Functions}
         -``\textit{Durable Functions is an extension of Azure
           Functions that lets you write stateful functions in a
           serverless compute environment. The extension lets you
           define stateful workflows by writing orchestrator functions
           and stateful entities by writing entity functions using the
           Azure Functions programming model.}''
       \item State machine, encoded into programming language, C\#,
         JS, Python, F\#, PowerShell
       \item API \alert{to run}
    \end{itemize}
\end{nrcanFrame}

\begin{nrcanFrame}
  \frametitle{FAAS Orchestration: GCP}
  \begin{itemize}
  \item \href{https://cloud.google.com/composer/}
         {Cloud Composer}
         -``\textit{A fully managed workflow orchestration service
           built on
           \href{https://airflow.apache.org/}{Apache Airflow}.}''
         \begin{itemize}
         \item State machine, Python-centric
         \item API to run and update?
         \end{itemize}
  \end{itemize}
\end{nrcanFrame}


%%%%%%%%%%%%%%%%%%%%%%%%%%%%%%%%%%%%%%%%%%%%%%%%%%%%%%%%%%%%%%%%%%%%%%%%%%%%%%%
%%%%%%%%% Last page (Canada word mark)
\nrcanLastFrame

\end{document}
